\section{Introduction}\label{sec:intro}

Measuring the stock of fixed physical capital is fundamental to macroeconomics, growth theory, and political economy. Capital stock estimates serve as the primary denominator for profit rates, capital--output ratios, capacity utilization metrics, and total factor productivity. Yet for Chile, long-run historical data on fixed capital have remained fragmented along three major lines: valuation (gross capacity versus net book value), asset coverage (residential versus non-residential construction and machinery), and temporal continuity across structural breaks in official national accounts.

\emph{This paper resolves this fragmentation by constructing a unified, stock--flow consistent (SFC) dataset of gross and net capital stocks, investment flows, retirements, depreciation, implicit depletion and depreciation rates, and capital price deflators for Chile spanning 125 years (1900--2024).} Building upon the foundational contributions of \citet{Hofman2000a}, \citet{DiazWagner2016}, \citet{TafunellDucoing2016}, \citet{PerezEyzaguirre2019}, and the Central Bank of Chile (\citealt{BCCH2024}), we introduce a integrated harmonization framework that resolves historical scale discrepancies, enforces stock--flow identity accounting, and extends constant-price 2003 Chilean Peso (CLP) series through 2024 while maintaining historical series continuity.

The theoretical necessity of this harmonization is operational rather than purely doctrinal. In empirical growth and political economy research, serious inferential errors arise when net book-value series are conflated with installed physical capacity. Net capital stocks ($K_n$) reflect accounting depreciation schedules, tax write-downs, and valuation adjustments; they measure surviving economic value and wealth. Gross capital stocks ($K_g$), by contrast, record physical installed capacity prior to scrapping; they provide the appropriate physical base for analyzing capital productivity, capacity utilization, and maximum output potential \citep{BasuRama2013, Shaikh2016}. By tracking both gross and net stock dynamics simultaneously within an explicit accounting structure, our framework allows researchers to examine how retirements (scrapping/depletion $R_t$) and consumption of fixed capital (depreciation $D_t$) diverge during structural crises, technological transformations, and industrialization phases.

Methodologically, our construction implements a stock--flow consistent law of motion that governs gross and net capital accumulation across five core asset categories: Machinery and Equipment (ME), Non-Residential Construction (NRC), Residential Construction (RC), Total Non-Residential Capital ($\text{NR} \equiv \text{ME} + \text{NRC}$), and Total Fixed Capital ($\text{Total} \equiv \text{ME} + \text{NRC} + \text{RC}$). The methodology enforces four strict principles:
\begin{enumerate}[label=(\roman*)]
    \item \textbf{Stock--Flow Identity}: Gross stock variations ($\Delta K_{g,t} = I_t - R_t$) and net stock variations ($\Delta K_{n,t} = I_t - D_t$) are dynamically linked through gross fixed capital formation ($I_t$).
    \item \textbf{Capacity Backlog Closure}: The difference between gross capacity and net value defines the capacity backlog ($B_t \equiv K_{g,t} - K_{n,t}$), whose motion obeys $\Delta B_t = R_t - D_t$.
    \item \textbf{Implicit Rate Identification}: Depreciation rates ($\delta_t \equiv D_t / K_{n,t-1}$) and depletion/scrappage rates ($z_t \equiv R_t / K_{g,t-1}$) are derived empirically from the joint motion of flows and stocks rather than imposed via arbitrary exogenous schedules.
    \item \textbf{Historical Continuity and Extension}: The dataset is built on a 2003 CLP historical baseline for 1900--1994 and extended to 2024 using Central Bank of Chile (BCCh) net stock and depreciation dynamics via asset-specific 1994 level-conversion ratios.
\end{enumerate}

The empirical contributions of this paper are fourfold. First, we deliver the first continuous, 125-year dataset reporting both gross ($K_g$) and net ($K_n$) capital stocks by detailed asset class for Chile (1900--2024) in unified constant 2003 CLP. Second, we provide a single, continuous capital price deflator ($P_k$, $2003=100$) from 1940 to 2024, splicing the official BCCh implicit deflator with chained ClioLab ratio growth. Third, we document long-run historical trajectories of implicit depreciation ($\delta_t$) and scrappage ($z_t$) rates, identifying major capital destruction and modernization waves during the 1930 Great Depression, the post-1973 structural transformation, and the 1982 debt crisis. Fourth, we provide a complete, AEA-compliant open-access Reproduction Package with machine-readable provenance, SHA-256 checksums, and a single-command build interface (\texttt{python run\_all.py}).

The remainder of this paper is organized as follows. Section~\ref{sec:literature} reviews historical national accounts and the relevant literature. Section~\ref{sec:framework} details the stock--flow accounting framework. Section~\ref{sec:harmonization} presents data sources, level normalization rules, asset composition allocations, and BCCh ratio calculation details. Section~\ref{sec:results} documents long-run patterns of Chilean accumulation from 1900 to 2024. Section~\ref{sec:validation} presents accounting validation tests and source reconciliation. Section~\ref{sec:conclusion} concludes and outlines macroeconomic applications.
